\section*{Chapter \ref{ch:g8:the-atmosphere} --- The Atmosphere}

\begin{solution}{exo:g8:the-atmosphere:1}
(a) \omterm{def:g8:the-atmosphere:layers}{stratosphere}; (b) \omterm{def:g8:the-atmosphere:layers}{troposphere}; (c) \omterm{def:g8:the-atmosphere:layers}{mesosphere}; (d) \omterm{def:g8:the-atmosphere:layers}{thermosphere}.
\end{solution}

\begin{solution}{exo:g8:the-atmosphere:2}
(a) \qty{15}{\degreeCelsius}. (b) About \qty{-56}{\degreeCelsius}.
(c) About \qty{-3}{\degreeCelsius}. (d) About \qty{-86}{\degreeCelsius}.
(e) From \qty{11}{km} to about \qty{20}{km} (and over a few kilometres
near \qty{50}{km} and just above \qty{86}{km}).
\end{solution}

\begin{solution}{exo:g8:the-atmosphere:3}
(a) False: it falls in the \omterm{def:g8:the-atmosphere:layers}{troposphere}, rises in the \omterm{def:g8:the-atmosphere:layers}{stratosphere}, falls
in the \omterm{def:g8:the-atmosphere:layers}{mesosphere} and rises again in the \omterm{def:g8:the-atmosphere:layers}{thermosphere}. (b) True.
(c) False: it is in the \omterm{def:g8:the-atmosphere:layers}{stratosphere}. (d) True.
\end{solution}

\begin{solution}{exo:g8:the-atmosphere:4}
(a) The \omterm{def:g8:the-atmosphere:layers}{troposphere}. (b) The \omterm{def:g8:the-atmosphere:layers}{stratosphere} (in the standard atmosphere the
\omterm{def:g8:the-atmosphere:layers}{tropopause} is at \qty{11}{km}). (c) The \omterm{def:g8:the-atmosphere:layers}{stratosphere}. (d) The \omterm{def:g8:the-atmosphere:layers}{mesosphere}.
\end{solution}

\begin{solution}{exo:g8:the-atmosphere:5}
(a) $20 - 6.5 \times 2 = 20 - 13 = \qty{7}{\degreeCelsius}$. (b) $20 -
6.5 \times 3 = 20 - 19.5 = \qty{0.5}{\degreeCelsius}$. (c) The air has to
cool by 20 degrees: $20 \div 6.5 \approx 3.1$, so at about
\qty{3.1}{km}.
\end{solution}

\begin{solution}{exo:g8:the-atmosphere:6}
(a) $3 - 0.5 = \qty{2.5}{km}$. (b) $6.5 \times 2.5 = 16.25$: about 16
degrees. (c) $18 - 16.25 \approx \qty{2}{\degreeCelsius}$.
\end{solution}

\begin{solution}{exo:g8:the-atmosphere:7}
(a) \qty{500}{hPa}, \qty{250}{hPa} and \qty{125}{hPa}. (b) Near the
ground the rule is good (\qty{500}{hPa} against \qty{506}{hPa}). Higher
up it gives too much: \qty{250}{hPa} against \qty{227}{hPa}, and
\qty{125}{hPa} against \qty{96}{hPa}. The halving rule is only a rough
guide.
\end{solution}

\begin{solution}{exo:g8:the-atmosphere:8}
(a) $10~\text{hPa} = 10^1$~hPa; $10^3 \div 10^1 = 10^2$: a hundred times
smaller. (b) $0.01~\text{hPa} = 10^{-2}$~hPa. (c) $10^3 \div 10^{-2} =
10^5$: a hundred thousand times smaller.
\end{solution}

\begin{solution}{exo:g8:the-atmosphere:9}
(a) $6 \div 1.5 = 4$ times wider. (b) As it rises, the air around it
presses less and less, so the gas inside spreads out and the balloon
grows until its skin breaks. (c) After the balloon bursts, the parachute
slows the fall of the radiosonde, so that it does not hurt anyone or
damage anything.
\end{solution}

\begin{solution}{exo:g8:the-atmosphere:10}
(a) $20 \div 7 \approx 2.9$ times thicker. (b) Closer to the \omterm{def:g3:continents-and-oceans:poles}{poles}:
\qty{11}{km} is \qty{4}{km} above \qty{7}{km} but \qty{9}{km} below \qty{20}{km} (it describes the middle \omterm{def:g6:latitude-longitude:latitude}{latitudes}, not the \omterm{def:g3:continents-and-oceans:equator}{equator}).
(c) At the start the \omterm{def:g8:the-atmosphere:layers}{tropopause} is at about \qty{7}{km}, below the
plane: the plane is in the \omterm{def:g8:the-atmosphere:layers}{stratosphere}. At the end the \omterm{def:g8:the-atmosphere:layers}{tropopause} is at
about \qty{20}{km}: the plane is in the \omterm{def:g8:the-atmosphere:layers}{troposphere}.
\end{solution}

\begin{solution}{exo:g8:the-atmosphere:11}
(a) $1 - 506 \div 1013 \approx 0.50$: \textbf{50~\%}. (b) $1 - 227 \div
1013 \approx 0.78$: \textbf{78~\%}. (c) $1 - 12 \div 1013 \approx 0.99$:
\textbf{99~\%}. (d) About 1~\%.
\end{solution}

\begin{solution}{exo:g8:the-atmosphere:12}
(a) $6371 \div 15 \approx 425$~km for each centimetre. (b) $100 \div 425
\approx 0.24$~cm, about \qty{2.4}{mm}. (c) $11 \div 425 \approx
0.026$~cm, about \qty{0.26}{mm}. (d) $0.26 \div 0.1 \approx 2.6$: two or
three sheets of paper.
\end{solution}

\begin{solution}{pb:g8:the-atmosphere:1}
\textbf{1.} $300 \times 10 = 3000$~m.

\textbf{2.} $6000 \div 300 = 20$ minutes; $\num{30000} \div 300 = 100$
minutes.

\textbf{3.} 100 minutes, that is 1~hour and 40~minutes.

\textbf{4.} At \qty{8}{km}: the \omterm{def:g8:the-atmosphere:layers}{troposphere} (the temperature is still
falling). At \qty{30}{km}: the \omterm{def:g8:the-atmosphere:layers}{stratosphere}.

\textbf{5.} $15 - (-50) = 65$ degrees.

\textbf{6.} $65 \div 10 = \qty{6.5}{\degreeCelsius}$ per kilometre.

\textbf{7.} From \qty{12}{km} to \qty{20}{km} it stays at
\qty{-56}{\degreeCelsius}. The cooling stops between \qty{10}{km} and
\qty{12}{km}: the \omterm{def:g8:the-atmosphere:layers}{tropopause} is at about \qty{11}{km}.

\textbf{8.} It rises: \qty{-55}{\degreeCelsius} at \qty{22}{km},
\qty{-47}{\degreeCelsius} at \qty{30}{km}. The balloon is in the
\omterm{def:g8:the-atmosphere:layers}{stratosphere}, warmed by the ozone that absorbs the Sun's ultraviolet
light.

\textbf{9.} $1013 \div 2 \approx \qty{507}{hPa}$; it is reached between
\qty{4}{km} (\qty{617}{hPa}) and \qty{6}{km} (\qty{472}{hPa}).

\textbf{10.} $1013 \div 4 \approx \qty{253}{hPa}$.

\textbf{11.} Between \qty{10}{km} (\qty{265}{hPa}) and \qty{12}{km}
(\qty{194}{hPa}).

\textbf{12.} From \qty{10}{km}, the pressure has to fall by $265 - 253 =
12$~hPa more, out of the 71~hPa it loses over 2~km: $12 \div 71 \times 2
\approx 0.34$~km. The pressure is a quarter of its sea-level value at
about \textbf{\qty{10.3}{km}}.
\end{solution}
